\section{What We Tried in Causal Training---and Why We Did Not Keep It}
\label{app:causal-design-history}

The retained method in Section~\ref{sec:causal-training} is the end of a much
larger search.  This appendix gives the useful lessons without requiring the
reader to follow the order in which kernels were written.  Internal version
names appear only here and in the provenance appendix.

\subsection{Projection precision and attention precision are different}

The first important separation is between learned projection layers and the
attention kernel.  QKV and output projections multiply model activations by
learned weights; attention then consumes the resulting Q, K, and V tensors.
Using FP4 successfully inside attention does not imply that learned weights
and activations can use the same format.

In an isolated 8B transformer block, E4M3 QKV projection had approximately
0.0375 relative-$L_2$ error, compared with 0.1459 for NVFP4 QKV.  When the
decoded low-precision operands were passed through exact bfloat16 attention,
the attention-output error was already 0.3614.  The native attention kernel added only
0.0216 relative-$L_2$ on those same represented operands.  In other words,
most error entered before the attention kernel.  This is why the main
experiments treat projection precision as a separate variable: E4M3 is the
numerical control, while NVFP4 is the higher-throughput arm.  This isolated
block does not by itself choose a final projection format.

An earlier fixed-head scaling rule also failed immediately on the 8B shape.
Replacing it with two-dimensional row-by-K16 Q/K scales reduced pre-clipping
gradient norms from roughly 20,000 to 161--184 in the short gate and restored
bfloat16-like behavior.  This was a representation fix, not a change to the
matrix schedule.

\subsection{Two independent backward scale failures}

Backward initially produced finite tensors that were nevertheless wrong.  Two
independent factors of 256 were involved:

\begin{enumerate}
  \item score reconstruction consumes a lifted statistic
  $\ell_i=8-L_i\log_2(e)$, where $L_i$ is the saved natural-log LSE.  Omitting
  $+8$ reconstructs the probability at $1/256$ of its intended scale;
  \item the gradient-output epilogue has its own $1/256$ correction.  It is a
  separate matrix-layout convention and must not be folded into the LSE fix.
\end{enumerate}

After those corrections, a checkpoint diagnostic exposed a different problem:
E4M3 rounded about 97.1\% of dO values to zero.  E5M2 has fewer fraction bits
but a wider exponent range; publishing dO in E5M2 reduced the observed zero
fraction to about 14.4\%.  The E5M2 publisher was only 0.7--0.9\% slower than
the E4M3 publisher in the matched gate.  This is why the final method uses two
different FP8 formats rather than treating all 8-bit values as interchangeable.

\subsection{How the native backward schedule evolved}

Table~\ref{tab:causal-version-map} maps internal development labels to the
idea each one tested.  The labels are not public method names.

Some measurements apply only to predecessor schedules.  In the v501 B2 owner4
route, dK and dV had unique writers, so only dQ required clearing at the exact
B2/S4096 shape.  A bounded profile of that same predecessor measured 21\% warp
occupancy, 30\% tensor activity, and 9\% external-memory activity despite 95\%
SM activity.  The retained v509 entry point always clears its outputs; no
equivalent performance profile exists for that final binary.  Reuse of rounded
P, dual-layout dS publication, and early tensor-memory release remain useful
design ideas, but predecessor counters are not evidence for final-route speed.

\begin{table}[htbp]
\centering
\caption{Internal backward version map.}
\label{tab:causal-version-map}
\small
\begin{tabularx}{\textwidth}{l p{.29\textwidth}X}
\toprule
Label & Main purpose & Outcome \\
\midrule
v416 & D64 native owner schedule & Like-for-like parity with the generated
native-exponential reference; used in early 1.2B integration. \\
v454/v482 & D128 B1/B2 ownership, rounded-P reuse, and early tensor-memory
release & 1.21--1.24$\times$ faster than the matched generated reference. \\
v501 & Corrected LSE lift, shape-specific clearing, and represented E4M3
gradient operands & Finite short-run systems prototype and basis of the
historical 8B bracket. \\
v503 & Common-row MXFP4 V approximation in backward & Faster than the first
MX attempt and tied end to end; the complete recipe failed its observed
distributed numerical gate, but the consumer was not isolated as the cause. \\
v506/v507 & Direct shared-MX producer and exact four-anchor consumer &
Numerically useful controls, but too slow for the production gate. \\
v509 & Exact forward NVFP4 score reconstruction with E5M2 dO & Retained
quantized causal-backward implementation; exact-batch B1/B2/B4 binaries are
validated for the Llama-style D128 shape. \\
\bottomrule
\end{tabularx}
\end{table}
\FloatBarrier

\subsection{Historical isolated backward timings}

The longest matched bfloat16 sequence sweep used a generated CuTe kernel, not
the retained implementation.  It nevertheless records how that earlier
low-precision path scaled from 512 to 16,384 tokens.

\begin{table}[htbp]
\centering
\caption{Historical isolated D64 causal backward on GB200.  The low-precision
route is a generated CuTe kernel; both columns include required output clears.}
\label{tab:d64-bf16-sweep}
\small
\begin{tabular}{rrrr}
\toprule
Sequence & Exact BF16 & Low precision & Speedup \\
\midrule
512   & 111.456 $\mu$s  & 102.176 $\mu$s  & 1.091$\times$ \\
1024  & 149.504 $\mu$s  & 138.560 $\mu$s  & 1.079$\times$ \\
2048  & 203.808 $\mu$s  & 189.632 $\mu$s  & 1.075$\times$ \\
4096  & 347.104 $\mu$s  & 319.808 $\mu$s  & 1.085$\times$ \\
8192  & 879.168 $\mu$s  & 770.848 $\mu$s  & 1.141$\times$ \\
16384 & 2765.984 $\mu$s & 2621.376 $\mu$s & 1.055$\times$ \\
\bottomrule
\end{tabular}
\end{table}
\FloatBarrier

At D128, the optimized native predecessor was compared with a generated
low-precision reference using the same represented operands.  This isolates
scheduling improvement; it is not a low-precision-versus-BF16 comparison.

\begin{table}[htbp]
\centering
\caption{Historical isolated D128 causal backward at S4096/Hq32/Hkv8 on
GB200.  The native column is the v454/v482 predecessor family.}
\label{tab:causal-backward-results}
\small
\begin{tabular}{lrrr}
\toprule
Shape & Generated reference & Native schedule & Speedup \\
\midrule
B1/D128 & 381.376 $\mu$s & 315.360 $\mu$s & 1.209$\times$ \\
B2/D128, rotation A & 620.032 $\mu$s & 514.272 $\mu$s & 1.206$\times$ \\
B2/D128, rotation B & 639.328 $\mu$s & 515.040 $\mu$s & 1.241$\times$ \\
\bottomrule
\end{tabular}
\end{table}
\FloatBarrier

\subsection{Why MXFP4 V in backward was not retained}

MXFP4 P/V is attractive in forward because the isolated D128/B2 core measures
259.072 $\mu$s, compared with 295.008 $\mu$s for FP8 P/V.  The full training
step, however, also needs V in a layout and format suitable for dP.  The first
MX route therefore published an MXFP4 V for forward and a second E4M3 V for
backward, consuming most of the isolated saving before backward began.

We tried to remove the second publication.  A shared-tile producer quantized
each D32-by-S32 V tile once and wrote both physical orientations.  An exact
four-anchor backward matched its represented oracle but took 581.632 $\mu$s,
versus 485.376 $\mu$s for the E4M3-V control.  The faster common-row consumer
reduced time but changed dQ/dK: cosines were about 0.989 and relative-$L_2$
about 0.145 against the E4M3-V route.  At whole-step scale, shared MX was only
0.194\% faster at the median and 0.102\% faster in sustained time---a tie, not
a useful throughput result.

Small S256 stochastic-rounding tests decoded MX values into E4M3 and therefore
served only as numerical proxies; they did not exercise a packed-MX kernel or
measure end-to-end time.  Those proxies and larger-batch checks did not change
the final choice.  The exact MX consumer was too slow, while the complete
approximate-consumer recipe failed the observed distributed numerical gate;
the failure was not isolated to that consumer.  The retained training path
therefore keeps E4M3 V views for backward and selects FP8 P/V in forward.

\subsection{Historical local speed brackets}

The early end-to-end brackets were valuable engineering checks, but they used
older learned-projection formats and cut-cross-entropy (CCE), so they are not
timings of the current dense-cross-entropy recipe.  Table
\ref{tab:historical-local-brackets} keeps them for context.

\begin{table}[htbp]
\centering
\caption{Historical saturated single-GPU brackets.  Each speedup is valid
within its row group but must not be transferred to the final training recipe.}
\label{tab:historical-local-brackets}
\small
\begin{tabular}{llrr}
\toprule
Model/shape & Attention route & Update time & Versus BF16 \\
\midrule
1.2B, B16/S4096 & BF16 FA4 & 673.396 ms & 1.000$\times$ \\
 & NVFP4-QK + FP8-PV & 615.682 ms & 1.094$\times$ \\
 & NVFP4-QK + MXFP4-PV & 614.842 ms & 1.095$\times$ \\
\midrule
8B, B2/S4096 & BF16 FA4 & 489.821 ms & 1.000$\times$ \\
 & NVFP4-QK + FP8-PV & 434.014 ms & 1.129$\times$ \\
 & NVFP4-QK + dual-published MXFP4-PV & 435.992 ms & 1.124$\times$ \\
\bottomrule
\end{tabular}
\end{table}
\FloatBarrier

The 8B FP8 and dual-published MX routes used the same backward.  Their backward
medians differed by only 0.084 ms, while the complete MX update was 1.978 ms
slower.  This explains the apparent paradox: an isolated MX forward kernel can
be materially faster even though its roughly one-millisecond saving per
32-layer model is small enough to be erased by publication and ordinary
step-level variation.

\subsection{Historical training context}

The September 1 snapshot contains two FP8-P/V trajectories that remain
non-divergent through a common 55.5-billion-token horizon.  Figure
\ref{fig:llama8b-training-curves} includes an older bfloat16 causal-FA4 run as
a sanity reference.  At that horizon, the one-billion-token exponential
moving-average losses are 2.534 for E4M3 projections with FP8 P/V, 2.542 for
NVFP4 projections with FP8 P/V, and 2.756 for the historical bfloat16 run.
These coordinates cannot be interpreted as relative model quality: the
bfloat16 run used a different topology, global batch, sample order, and
history-sampling density.

\begin{figure}[H]
\centering
\reportplot[0.27\textheight]{figures/llama8b_training_curves.pdf}
\caption{Historical 8-billion-parameter training context through the common
55.5-billion-token horizon.  Thin lines show downsampled training loss and
thick lines show a causal exponential moving average with a one-billion-token
half-life.  The FP8-P/V arms use 64-way replicated data parallelism with
global batch 64.  The bfloat16 reference uses 32-way fully sharded data
parallelism with global batch 32; it is not a paired control.}
\label{fig:llama8b-training-curves}
\end{figure}

\subsection{What the distributed failures established}

The first long distributed MX run combined learned NVFP4 projections, CCE,
and the common-row MX backward, so its failure did not identify a single
cause.  At a matched checkpoint it had loss 7.0805 versus 6.5887 for FP8 and a
pre-clipping gradient norm of 157,696 versus 6.41.  It later exceeded 1.5
million despite gradient clipping.

The final four-arm diagnostic removed that ambiguity by crossing learned
projection format with forward P/V format while holding the backward fixed.
All four curves remain close through update 300 (78.6 million tokens).  The
MXFP4 arms visibly separate near update 400 (104.9 million tokens) and are
unambiguously split by update 500 (131.1 million tokens): loss is 7.97 for
both MXFP4 arms, compared with 5.41 for their FP8 controls.  Both MX routes
continue to fail under different learned-projection formats, whereas both FP8
routes remain non-divergent through the snapshot.  Thus forward MXFP4 P/V, or
the model state it induces, is the common separator.  This is strong
factorial evidence, but not a formal proof that one kernel instruction causes
divergence.

\begin{table}[htbp]
\centering
\caption{Initial frozen rolling-log cutoff, retained for provenance.  Jobs
began at different times, so these are status observations rather than
aligned loss or throughput comparisons.}
\label{tab:v509-four-arm-cutoff}
\small
\begin{tabular}{llrrrrl}
\toprule
Projections & Forward P/V & Update & Tokens & Loss & Grad norm & Status \\
\midrule
E4M3  & FP8   & 10,381 & 2.721B & 2.9743 & 0.2197 & working at cutoff \\
E4M3  & MXFP4 & 10,075 & 2.641B & 8.8265 & 358,400 & diverged \\
NVFP4 & FP8   & 10,884 & 2.853B & 3.1637 & 0.2051 & working at cutoff \\
NVFP4 & MXFP4 & 11,061 & 2.900B & 8.4801 & 3,915,776 & diverged \\
\bottomrule
\end{tabular}
\end{table}
\FloatBarrier

The old rolling-log receipt first retained bad observations near update 4020
(1.054B tokens) for E4M3+MX and update 4722 (1.238B tokens) for NVFP4+MX.
The complete histories in Figure~\ref{fig:llama8b-mxfp4-divergence} show that
these were truncation markers, not onset estimates.  At the common
55.5-billion-token horizon, the one-billion-token moving-average losses are
8.94 and 8.18 for E4M3+MX and NVFP4+MX, compared with 2.53 and 2.54 for their
FP8 controls; the MX gradient norms repeatedly reach the millions.  The
receipt hashes and scientific route identities are listed in
Appendix~\ref{app:causal-training-provenance}.

A later B4 campaign supported the same route selection at a larger effective
batch.  The NVFP4-projection/MXFP4-PV arm diverged shortly after its initially
matched region.  A second arm changed the learned projections to E4M3 while
retaining MXFP4 P/V; it also developed repeated, finite gradient explosions
and was cancelled at update 2550.  Its last completed held-out validation was
7.4481 at update 2384.  The matched bfloat16 and
NVFP4-projection/FP8-PV trajectories remained stable.  The two failed runs
change projection precision but retain MXFP4 P/V, making the P/V
representation the common factor.  This does not prove that MXFP4 is
fundamentally unsuitable or show whether forward quantization, the saved V
payload, or its backward use is responsible.  We therefore show the current
B4 failure separately in Figure~\ref{fig:llama8b-b4-mxfp4-failure} rather than
mixing it into the healthy matched-training plot.

\begin{figure}[p]
\centering
\reportplot[0.34\textheight]{figures/llama8b_b4_mxfp4_failure.pdf}
\caption{Matched B4 MXFP4-P/V failure diagnostic.  This arm uses E4M3 learned
projections, NVFP4 Q/K inside attention, and MXFP4 P/V.  The top panel shows
training and held-out validation loss; the lower panel shows the pre-clipping
gradient norm.  The run was cancelled after update 2550 and is excluded from
the healthy-route throughput comparison.}
\label{fig:llama8b-b4-mxfp4-failure}
\end{figure}

\begin{figure}[p]
\centering
\reportplot[0.34\textheight]{figures/llama8b_mxfp4_divergence.pdf}
\caption{Separated diagnostic for the two diverging MXFP4-P/V arms.  The top
panel shows loss; the lower panel shows the maximum pre-clipping gradient norm
within each 100-update display bin on a logarithmic scale.  The two FP8
control arms are muted and use the corresponding projection formats.  The
inset expands the first 0.55 billion tokens and qualitatively marks the
interval between approximately 0.08 and 0.14 billion tokens in which the
curves first separate.}
\label{fig:llama8b-mxfp4-divergence}
\end{figure}
\FloatBarrier

\subsection{Evidence boundary}

The local brackets measure the retained backward, the complete attention
sublayer, and a full 8B update.  The distributed histories show that the FP8
route remains non-divergent and continues descending over the observed window,
while the MXFP4 route does not.  The same-recipe B4 bfloat16 control provides
same-update held-out validation and throughput over the completed
100-billion-token schedule.
Because there is one trajectory per route, these data do not support
run-to-run variability or statistical-equivalence claims.  Input stalls and
checkpoint windows remain in the reported throughput distribution.
